{\large\textbf{Glass ball}}

The first photo here is taken with a digital camera and shows a glass ball,
backlit with diffuse dichromatic light which has only two narrow spectral lines
(red 630\,nm and violet 400\,nm). This diffuse light comes from the white floor
(marked with `1' in the figure) and the white walls (marked with `2'), both
illuminated with violet and red LED lamps. The sensor of the camera has only
red, blue and green sensors so that violet light appears in the photo as blue.
The photo is taken from a distance much greater than the radius of the ball. On
the back side of the ball, a very narrow opaque thread is glued to the glass
surface, forming an arc of a great circle on the ball. In the photo, the thread
is obstructed by the ball and cannot be seen directly. However, hugely deformed
images of a very short segment of the thread are seen as blue (marked with
`b') and red (`r') ellipses. The letter `p' indicates purple-coloured areas in
the photo.

\begin{center}\includegraphics[width=0.42\linewidth]{figures/EuPhO_2021_Q3_fig1.jpg}\end{center}

In the first photo, the centre of the ball is marked with a cross, and the
perimeter of the ball is traced with a dashed line. You can find a larger
version of the first photo on a separate sheet. You may take distance
measurements there. In the larger photo, the boundary between the red and
purple regions is also traced with a dashed line.

The second photo is taken while illuminating with a white LED with the ball
turned so that the thread can be seen directly.

\begin{center}\includegraphics[width=0.42\linewidth]{figures/EuPhO_2021_Q3_fig2.jpg}\end{center}

a) Explain qualitatively using a ray diagram why a segment of the thread is
seen as a closed loop in the first photo.

b) Determine the coefficient of refraction $n_R$ for the red light.

c) Determine the difference of the coefficients of refraction
$\Delta n \equiv n_V - n_R$ for the red and violet light (with $n_V$ denoting
the coefficient of refraction for the violet light)

% the separate sheet: larger version of the first photo (page 2 of the source)
\begin{center}\includegraphics[width=1.00\linewidth]{figures/EuPhO_2021_Q3_fig3.jpg}\end{center}
